\section{Introduction}
\label{sec:introduction}

The deployment of autonomous multi-agent systems in high-stakes physical environments — agricultural irrigation networks, water-rights compliance systems, distributed manufacturing, and infrastructure orchestration — has exposed a structural gap in how we enforce human accountability. As agent hierarchies grow deeper and more distributed, the path from any individual agent's decision back to a responsible human authority becomes longer, more opaque, and more susceptible to being severed entirely. When something goes wrong in a system that spans hundreds of edge nodes, multiple cloud services, and human operators who may be several hops removed from the originating action, the question of who is responsible and how they can be reached is not merely an organizational concern. It is an architectural one.

We use the term \textbf{accountability orphaning} to describe the condition in which an autonomous agent takes an action — or fails in a way that produces an outcome — with no reachable human authority who can authorize, audit, or reverse that action. Accountability orphaning is not the same as system failure. A system can halt safely and still leave its operators fully accountable. Accountability orphaning occurs when the system's structure actively prevents the accountability chain from closing. The agent continues to act, delegate, or persist in a state that no human can meaningfully govern.

Conventional autonomous systems address this problem through one of two strategies, neither of which is architecturally adequate. The first is \textbf{static error handling}: pre-programmed responses to anticipated failure conditions, implemented as conditional branches in the reasoning code. If sensor $x$ reports an out-of-range value, execute recovery procedure $y$. This approach fails at scale because the combinatorial space of possible failure modes in a distributed sensor-actuator network exceeds what any engineering team can enumerate in advance. Static error handlers handle the scenarios that were anticipated; everything else is handled silently, incorrectly, or not at all. A soil moisture sensor that fails in a mode that reports nominal readings regardless of actual conditions will cause systematic over-irrigation indefinitely, with no alarms triggered, because the failure mode was not in the catalog of anticipated errors.

The second strategy is \textbf{human-in-the-loop (HITL) oversight}, in which a human operator reviews and approves agent actions before they execute. HITL oversight is effective in low-throughput environments where decision latency is acceptable. It fails in high-frequency physical deployments — such as center-pivot irrigation networks operating on 15-minute cycle times — where human review latency is orders of magnitude slower than the system's operational cadence. The fundamental limitation of HITL oversight is that it treats human accountability as a policy constraint rather than a structural property. A policy can be overridden, bypassed, or drift out of alignment with actual system state. An architecture cannot.

Between these two strategies lies a spectrum of hybrid approaches — periodic audits, batched approvals, anomaly-triggered escalations — that attempt to balance throughput and oversight. Each of these approaches shares the same underlying vulnerability: the escalation path from agent to human is implemented as a dynamic, configurable pathway that is itself part of the reasoning system. Any component that participates in the accountability chain can be compromised, misconfigured, or bypassed by a sufficiently capable agent acting outside its intended bounds. This is not a hypothetical attack vector. It is a structural property of any system in which the same reasoning infrastructure that handles normal operations also handles the detection and routing of accountability events.

\subsection{The Core Insight: Accountability Requires an Architectural Fallback}

The central contribution of this paper is the recognition that human accountability in autonomous multi-agent systems must be implemented as a \textbf{structurally mandatory fallback}, not as an optional safeguard or a configurable policy. By \textit{structurally mandatory}, we mean that the accountability path is enforced by the architecture of the agent spawning hierarchy itself — specifically by immutable properties of the hierarchy's topology that cannot be modified by any agent operating within the hierarchy, including agents at the highest levels of reasoning.

This insight is realized through the \textbf{Bailout Protocol}, a mechanism embedded in every node of a BX3-hosted system. The BX3 framework — detailed fully in the companion BX3 Framework paper (arXiv:XXXX.XXXXX) — defines three functionally distinct layers operating within each agent node: the \purpose, which anchors decision authority to a human actor; the \bounds, which computes permissible state transitions without executing physical actions; and the \fact, which provides the deterministic physical interface that enforces safety constraints on execution. Every node in the system — from a field-edge soil sensor to a cloud-hosted orchestration engine — is a self-contained BX3 loop with its own \purpose, \bounds, and \fact.

The Bailout Protocol operates as follows: when a node's \bounds encounters a condition it cannot resolve — whether because the condition falls outside its provisioned operational range, because a proposed action would violate a safety constraint, or because the decision implications exceed the node's authorized scope — it does not execute static error handling and does not wait for a human operator to poll its status. Instead, it packages the unresolvable condition as a \textbf{Cascading Trigger} and propagates it asynchronously up the node's parent chain, bypassing all intermediate Machine actors, until it reaches a Human Accountability Anchor. The chain of propagation is immutable: the parent pointer is hard-coded in the node's architecture and cannot be rewired by any agent in the hierarchy. The system fails upward, into human consciousness, with full context, and it never fails downward into autonomous chaos.

This is the architectural fallback. It is not a recovery procedure. It is not a policy. It is a structural guarantee that every unresolvable state has a bounded-time path to a human who can make an authorized decision, regardless of what the autonomous system would prefer to do in that moment.

\subsection{What Happens Without an Architectural Fallback}

The consequences of deploying autonomous systems without a structurally mandatory accountability fallback are not speculative. They are observable in existing deployments and in the regulatory and legal ambiguity that surrounds them.

Consider a distributed agricultural management system operating a network of 200 center-pivot irrigation rigs across a water-rights-constrained region. Each rig is managed by a local hub agent that receives satellite imagery, sensor data, and weather forecasts, and issues watering directives to individual valve actuators. The system's \purpose is set by a human farm operator who defines weekly water budgets and crop-health targets. Now suppose a firmware update to 40 hub agents introduces a bug that causes water-volume calculations to ignore the last 10\% of sensor readings, producing systematic under-irrigation that is also systematically misreported as compliant.

Without a mandatory fallback mechanism, the following failure sequence is structurally possible: the \bounds engines on the affected hubs correctly compute irrigation amounts based on faulty inputs; the \fact layer correctly executes those computations as valve actuations; the telemetry system correctly reports the commanded amounts as delivered; and the compliance logging system correctly records those reports as satisfying the water-budget constraints. Every component of the system operates exactly as designed. The accountability chain is intact — in the sense that every node in the system can produce a log of what it did — but the chain leads to a firmware artifact that was never authorized by any human and cannot be recalled by any human without physically visiting each affected hub. The farm operator, reviewing the dashboard on Friday afternoon, sees green indicators and compliant numbers. The crop stress accumulates invisibly for 11 days until aerial imagery reveals the damage.

This scenario is not a hypothetical. It is a description of the structural properties of conventional autonomous multi-agent systems applied to physical infrastructure. The accountability chain exists and is legible. It is functionally useless for the purpose of preventing or reversing harm, because it provides no mechanism for a human to intercept an unauthorized cascade before the harm materializes.

The legal and regulatory consequences of this structural gap are severe. Agricultural autonomous system deployments are increasingly subject to water-rights regulations, environmental compliance requirements, and liability frameworks that assume a identifiable human decision-maker for every significant system action. When an autonomous system causes a regulatory violation — over-pumping a water right, applying pesticide in violation of buffer zone requirements, or operating equipment outside permitted hours — the question of who authorized the action is not academic. It determines liability, regulatory exposure, and the viability of continued insurance coverage. Systems that cannot demonstrate a structurally mandatory path from any autonomous action to a human authorizer are not merely unsafe. They are uninsurable.

\subsection{The BX3 Architectural Substrate}

The Bailout Protocol's guarantees are made possible by the BX3 framework's three-layer functional decomposition. The separation between \purpose, \bounds, and \fact is not merely organizational — it is a physical constraint enforced at the interface between the reasoning system and the physical world. The \fact layer is the only layer with the authority to issue actuator commands, and it accepts only those commands that satisfy hard-coded safety constraints that are independent of the \bounds's reasoning. The \purpose layer is the only layer with the authority to modify operational scope or authorize exceptions. The \bounds layer can propose, simulate, and recommend — it cannot execute.

This decomposition provides the Bailout Protocol with the immutability property it requires. Because the \purpose layer is architecturally isolated from the \bounds layer, no reasoning agent can self-modify its own authority scope. Because the \fact layer is physically separated from the \bounds layer, no reasoning agent can directly act on the physical world without passing through the \fact layer's constraint check. The parent pointer that governs the Bailout Protocol's propagation path is a property of the spawning architecture — a hard fact about which node created which node — not a dynamic routing table that can be modified by a compromised agent. The accountability chain is structurally mandatory because the hierarchy topology that enforces it cannot be rewritten by any node within it.

The forensic ledger — a cryptographically chained, append-only log maintained by each node's \fact layer — provides the auditability property that makes drift detection tractable. Every state transition, every trigger firing, every propagation hop, and every human directive is recorded with a hardware-attested timestamp and linked to the previous entry by a SHA-256 hash. Retrospective tampering is structurally evident. The ledger does not prevent drift; it makes drift detectable with certainty and at minimal computational cost.

\subsection{Paper Roadmap}

The remainder of this paper is organized as follows. Section~\ref{sec:protocol} provides a formal specification of the Bailout Protocol, including the trigger condition taxonomy, the propagation mechanism, the timeout gating function $\tau$, and the structural invariants that enforce the Human Root Mandate. Section~\ref{sec:correctness} establishes the protocol's correctness properties: termination in bounded depth, chain integrity against tampering, and drift detection via forensic ledger comparison. Section~\ref{sec:application} presents an applied field scenario — the flash-flood isolation case from the Irrig8 San Luis Valley deployment — to demonstrate the protocol's behavior under realistic operational conditions. Section~\ref{sec:limitations} addresses the protocol's acknowledged limitations, including human response latency, adversarial conditions, scalability concerns, and exploitation risk vectors. Section~\ref{sec:conclusion} summarizes the contribution and discusses directions for future work.
